\begin{solapartat}
\punts{0,6} Atès que el període d'oscil·lació depèn de la massa i de la constant elàstica:
\[ T = 2\pi\sqrt{\frac{m}{k}} \;\Rightarrow\; m = k\left(\frac{T}{2\pi}\right)^{2} \]
A partir de la mesura del període i coneguda la constant elàstica podem determinar la massa d'un cos.

Encara més, si en lloc de comptar el temps per fer una oscil·lació complerta mesurem el temps que triga a fer un cert nombre d'oscil·lacions millorem la precisió amb què determinem el període i d'aquesta manera podem determinar amb més precisió la massa de l'objecte.

\punts{0,65} Com que el període només depèn de la massa i la constant elàstica, no depèn de la intensitat del camp gravitatori, la mesura es la mateixa a la Terra que a la Lluna.

De fet, a la Lluna serà una mica més precisa atès que no tenim l'efecte de la fricció amb l'aire.
\end{solapartat}

\begin{solapartat}
\punts{0,1}
\[ A = 0,06\ \mathrm{m} \]
\[ \omega = 2\pi f = 20\pi\ \mathrm{rad/s} \]

\destaca{Primera opció}

\begin{itemize}[leftmargin=2em,label={}]
\item \punts{0,1} Equació del MHS:
\[ x(t) = A\sin(\omega t + \varphi_0) \]

\item \punts{0,2}
\[ v(t) = \frac{dx(t)}{dt} = A\omega\cos(\omega t + \varphi_0) \]

\item \punts{0,1} llavors la velocitat màxima és:
\[ v_{\text{màx}} = A\omega = 3,77\ \mathrm{m/s} \]

\item \punts{0,2}
\[ a(t) = \frac{dv(t)}{dt} = -A\omega^2\sin(\omega t + \varphi_0) = -x\omega^2 \]

\item \punts{0,1} llavors l'acceleració màxima és:
\[ a_{\text{màx}} = A\omega^2 = 237\ \mathrm{m/s^2} \]

\item \punts{0,2} Si a $t=0$ l'acceleració és màxima, llavors $x$ és mínima,
\[ x(t=0) = -A \]
\[ -A = A\sin(\varphi_0) \;\Rightarrow\; \varphi_0 = \mathrm{ArcSin}(-1) = -\frac{\pi}{2}\ \mathrm{rad}. \]

\item \punts{0,25} Finalment, l'equació del moviment és:
\[ x(t) = 0,06\sin(20\pi t - \pi/2),\ \text{$x$ en m i $t$ en s.} \]
\end{itemize}

\destaca{Alternativament:}

\begin{itemize}[leftmargin=2em,label={}]
\item \punts{0,1} Equació del MHS:
\[ x(t) = A\cos(\omega t + \varphi_0) \]

\item \punts{0,2}
\[ v(t) = \frac{dx(t)}{dt} = -A\omega\sin(\omega t + \varphi_0) \]

\item \punts{0,1} llavors la velocitat màxima és:
\[ v_{\text{màx}} = A\omega = 3,77\ \mathrm{m/s} \]

\item \punts{0,2}
\[ a(t) = \frac{dv(t)}{dt} = -A\omega^2\cos(\omega t + \varphi_0) = -x\omega^2 \]

\item \punts{0,1} llavors l'acceleració màxima és:
\[ a_{\text{màx}} = A\omega^2 = 237\ \mathrm{m/s^2} \]

\item \punts{0,2} Si a $t=0$ l'acceleració és màxima, llavors $x$ és mínima,
\[ x(t=0) = -A \]
\[ -A = A\cos(\varphi_0) \;\Rightarrow\; \varphi_0 = \mathrm{ArcCos}(-1) = \pi\ \mathrm{rad}. \]

\item \punts{0,25} Finalment, l'equació del moviment és:
\[ x(t) = 0,06\cos(20\pi t + \pi),\ \text{$x$ en m i $t$ en s.} \]
\end{itemize}
\end{solapartat}
